
%(BEGIN_QUESTION)
% Copyright 2006, Tony R. Kuphaldt, released under the Creative Commons Attribution License (v 1.0)
% This means you may do almost anything with this work of mine, so long as you give me proper credit

Når hydrogen og oksygen kombineres under tilstrekkelig trykk og/eller temperatur, reagerer de og danner vann (H$_{2}$O) i et enkelt eksempel på {\it forbrenning}. Hvis dette er de eneste grunnstoffene som deltar, og prøvene er rene, er vann eneste reaksjonsprodukt (i tillegg til varmeenergi).

Hydrogen forekommer som gass ved romtemperatur og atmosfærisk trykk som atompar (et {\it diatomisk} molekyl), skrevet H$_{2}$. Oksygen forekommer tilsvarende som O$_{2}$.

Hvis jeg har en mengde rent hydrogengass på 250 mol, hvor mange mol rent oksygengass trengs for å reagere fullstendig med (forbrenne) hydrogenprøven? Anta gasstrykk 830 mm HgA.

\vskip 20pt \vbox{\hrule \hbox{\strut \vrule{} {\bf Forslag til sokratisk diskusjon} \vrule} \hrule}

\begin{itemize}
\item{} Bør $\Delta H$-verdien for forbrenningsreaksjonen være positiv eller negativ?
\item{} Hovedmotorene på Space Shuttle bruker hydrogen som drivstoff og oksygen som oksidator. Dermed er vanndamp hovedutslippet. Betyr det at motorene er helt ikke-forurensende?
\end{itemize}

\underbar{file i00569}
%(END_QUESTION)





%(BEGIN_ANSWER)

125 mol oksygengass kreves for å forbrenne 250 mol hydrogengass fullstendig, fordi antall hydrogen- og oksygenatomer per mol er like, og ett oksygenatom trengs per to hydrogenatomer.

\vskip 10pt

Trykket 830 mm HgA er overflødig informasjon, lagt inn for å utfordre deg til å vurdere om informasjonen er relevant for løsningen.

\vskip 10pt

Spørsmålet om rakettmotorene er forurensende krever vurdering utover selve forbrenningen i motoren. Selv om motoren hovedsakelig slipper ut vanndamp (pluss litt ureaktert hydrogen og oksygen), er prosessen for å produsere hydrogenbrensel i seg selv forurensende. Store mengder rent hydrogen produseres ofte ved {\it reforming}, der hydrokarboner som naturgass (CH$_{4}$) omdannes til hydrogen (H$_{2}$) og karbondioksid (CO$_{2}$), hvor sistnevnte er en klimagass.

%(END_ANSWER)





%(BEGIN_NOTES)


%INDEX% Kjemi, støkiometri: reaksjonsmengder

%(END_NOTES)

